\chapter{Background}
\label{background}

\section{Important results}

Start with some notation and conventions that will be used throughout the thesis.

We will now present some important results that will be used later in the thesis.

\begin{theorem}[Bertini]\label{thm:bertini}
    If $\mathcal{D}$ is a linear system on a variety $X$ in characteristic $0$, then the general member of $\mathcal{D}$ is smooth away from the base locus of $\mathcal{D}$ and the singular locus of $X$.
\end{theorem}
\todo[inline]{Add notation (this is from EH): A $\textit{linear system}$ on a scheme $X$ is a pair $\mathcal{D} = (\mathcal{L}, V)$, where $\mathcal{L}$ is a line bundle on $X$ and $V \subseteq H^0(X, \mathcal{L})$ a vector space of sections.}

We refer to (!?) for a proof. \cite{EH16}
\todo[inline]{Find a good reference for this theorem.}

% \begin{theorem}[Bezout]\label{thm:bezout}
%     Let $C_1, C_2 \subseteq \mathbb{P}^2$ be two projective plane curves of degree $d_1$ and $d_2$ respectively, with no common components. Then $C_1$ and $C_2$ intersect in exactly $d_1 d_2$ points, counted with multiplicity.
% \end{theorem}
\begin{theorem}[Bézout]\label{thm:bezout}
    Let $X, Y \subseteq \mathbb{P}^n$ be two closed subvarieties having complementary dimension, i.e. $\dim(X) + \dim(Y) = n$. If $X$ and $Y$ intersect transversely, then they will intersect in exactly $\deg(X) \cdot \deg(Y)$ points. Furthermore, if the intersection is finite, then the number of intersection points (counted with multiplicity) is at most $\deg(X) \cdot \deg(Y)$.
\end{theorem}
We refer to (!?) for a proof. \cite{EH16} corollary 2.4.
\todo[inline]{Find a good reference for this theorem. Don't require transverse intersections.}
\todo[inline]{If X and Y intersect in finitely many points, then the number of intersection points is at most deg(x)deg(Y).}


%-----------------------------------------------------------------------------
% \begin{definition}[\cite{Ha77}, 12.7]
%     Let $Y$ be a topological space. A function $\phi\colon Y \to \mathbb{Z}$ is \emph{upper semicontinuous} if for each $y \in Y$, there is an open neighborhood $U$ of $y$, such that for all $y' \in U, \phi(y') \leq \phi(y)$ 
% \end{definition}

% \begin{lemma}[\cite{EO25}, Cor. 10.59]
%     Let $k$ be a field, and let $f \colon X \to Y$ be a morphism of projective schemes over $k$ Then $f$ is proper. In particular, the image $f(X) \subseteq Y$ is closed in $Y$.
% \end{lemma}

\begin{theorem}[\cite{Vak25} 11.4.1]
    Suppose $\pi \colon X \to Y$ is a morphism of irreducible $k$-varieties, with $\dim X = m$ and $\dim Y = n$. Then there exists a nonempty open subset $U \subseteq Y$ such that for all $q \in U$, the fiber over $q$ has pure dimension $m-n$, or is empty.
\end{theorem}

\begin{lemma}\label{lem:irred_fiber}
    Let $\pi \colon X \to Y$ be a proper morphism to an irreducible variety, and all the fibers are nonempty and irreducible of the same dimension. Then $X$ is irreducible.
\end{lemma}
This can be found in \cite{EO25}, exercise 14.2.21.
\todo[inline]{Sjekk at oppgaven er samme plass.}

\begin{theorem}[Dimension formula for fibers] \label{thm:fiber}
    Let $k$ be a field and let $f\colon X \to Y$ be a dominant morphism of irreducible projective varieties over $k$. Let $\eta$ denote the generic point of $Y$ and write $X_\eta = X \times_Y \mathrm{Spec}\,k(Y)$ for the generic fiber. Then
    \[
        \dim X \;=\; \dim Y \;+\; \dim X_\eta.
    \]
\end{theorem}
\todo[inline]{Write the proof, and/or find reference. Maybe use \cite{Vak25} 11.4.2.}
% \begin{proof}
%     For irreducible varieties over a field, the Krull dimension equals the transcendence degree of the function field over $k$. Dominance of $f$ gives an inclusion of function fields $k(Y)\hookrightarrow k(X)$ and $k(X)$ is a finitely generated field extension of $k(Y)$. Hence
%     \[
%         \dim X = \operatorname{trdeg}_k k(X) = \operatorname{trdeg}_k k(Y) + \operatorname{trdeg}_{k(Y)} k(X).
%     \]
%     By definition $\operatorname{trdeg}_{k(Y)} k(X)=\dim X_\eta$, and $\operatorname{trdeg}_k k(Y)=\dim Y$, which yields the claimed equality.
% \end{proof}
% \begin{definition}[\cite{Ha77}, 12.7]
%     Let $Y$ be a topological space. A function $\phi\colon Y \to \mathbb{Z}$ is \emph{upper semicontinuous} if for each $y \in Y$, there is an open neighborhood $U$ of $y$, such that for all $y' \in U, \phi(y') \leq \phi(y)$ 
% \end{definition}

% \begin{lemma}[\cite{EO25}, Cor. 10.59]
%     Let $k$ be a field, and let $f \colon X \to Y$ be a morphism of projective schemes over $k$ Then $f$ is proper. In particular, the image $f(X) \subseteq Y$ is closed in $Y$.
% \end{lemma}

\begin{theorem}[Dimension formula for fibers] \label{thm:fiber}
    Let $k$ be a field and let $f\colon X \to Y$ be a dominant morphism of irreducible projective varieties over $k$. Let $\eta$ denote the generic point of $Y$ and write $X_\eta = X \times_Y \mathrm{Spec}\,k(Y)$ for the generic fibre. Then
    \[
        \dim X \;=\; \dim Y \;+\; \dim X_\eta.
    \]
\end{theorem}
\todo[inline]{Write the proof, and/or find reference. Maybe use Vakil 12.4.2.}
% \begin{proof}
%     For irreducible varieties over a field, the Krull dimension equals the transcendence degree of the function field over $k$. Dominance of $f$ gives an inclusion of function fields $k(Y)\hookrightarrow k(X)$ and $k(X)$ is a finitely generated field extension of $k(Y)$. Hence
%     \[
%         \dim X = \operatorname{trdeg}_k k(X) = \operatorname{trdeg}_k k(Y) + \operatorname{trdeg}_{k(Y)} k(X).
%     \]
%     By definition $\operatorname{trdeg}_{k(Y)} k(X)=\dim X_\eta$, and $\operatorname{trdeg}_k k(Y)=\dim Y$, which yields the claimed equality.
% \end{proof}